\begin{tikzpicture}[
  tramo/.style={draw=ESCOMGris, dotted, line width=0.5pt},
  tramoL/.style={font=\diag\bfseries, text=ESCOMAzulOscuro, fill=white,
    anchor=west, inner sep=1.5pt}]
  \Actores{op/1.5/{Operador\\ (Qadra)},api/5.6/{Servicio de la API},ev/9.9/{Catálogo y evaluador},pg/13.8/{PostgreSQL}}{-10.1}

  % ==== Preparar ====
  \draw[tramo] (0,-0.75) -- (16,-0.75);
  \node[tramoL] at (0,-0.75) {Preparar: revisión prospectiva};
  \Msg{1.5}{5.6}{-1.3}{\textbf{1} resultado + plazo en una instrucción}
  \Paso{5.6}{-1.42}{4.6}{Autentica; exige gestionar resultados y plazos (Owner o Litigante asignado)}
  \Msg{5.6}{13.8}{-2.45}{resuelve fuentes y soporte existentes}
  \Paso{5.6}{-2.57}{4.6}{Descifra y valida la versión exacta del soporte fuera de la transacción de escritura}
  \Msg{5.6}{9.9}{-3.6}{perfil, calendario y valores propuestos}
  \Paso{9.9}{-3.72}{5.4}{Vía prospectiva: el resultado aún no tiene revisión persistida; el reloj observado no es el instante definitivo}
  \Msg[msgR]{9.9}{5.6}{-4.72}{cálculo o bloqueo explicado}
  \Paso{5.6}{-4.84}{4.6}{Reautentica. Huella \textbf{HRDL1} de la instrucción; sin evento, recibo final ni reserva}
  \Msg[msgR]{5.6}{1.5}{-5.77}{instrucción revisable}

  % ==== Confirmar ====
  \draw[tramo] (0,-6.07) -- (16,-6.07);
  \node[tramoL] at (0,-6.07) {Confirmar: una sola transacción};
  \Msg{1.5}{5.6}{-6.52}{\textbf{2} confirma HRDL1}
  \Paso{5.6}{-6.64}{2.2}{Reautentica}
  \Msg{5.6}{13.8}{-7.32}{transacción}
  \Paso[pasoOsc]{13.0}{-7.44}{5.6}{Bloqueo común; reautoriza; vuelve a resolver y compara los registros cifrados; inserta el resultado R1 y obtiene su evento real; prepara el plazo R1 por la vía ordinaria; exige el cálculo revisado; guarda el origen con \textbf{HRDC1} y la auditoría}
  \Paso[pasoError]{13.0}{-8.95}{5.6}{Revierte todo si falla el plazo, el origen o la auditoría}
  \Msg[msgR]{13.8}{5.6}{-9.6}{confirmado}
  \Msg[msgR]{5.6}{1.5}{-9.95}{dos detalles R1 y el origen conjunto}

  % ==== Recuperar ====
  \draw[tramo] (0,-10.25) -- (16,-10.25);
  \node[tramoL] at (0,-10.25) {Respuesta perdida};
  \node[font=\diag, text=ESCOMTinta, align=left, text width=15.6cm, anchor=north west] at (0,-10.5)
    {Conciliación como en la \cref{fig:impl-audiencias-recibo}, buscando el origen conjunto: una confirmación idéntica concurrente recupera el mismo origen, una distinta entra en conflicto y nunca hay reenvío automático.};
\end{tikzpicture}
